% Boxed example for Section 2.3 (Machine-actionable guidance for autonomous
% data retrieval workflows). Requires \usepackage{tcolorbox} in the preamble.
% Drop % Boxed example for Section 2.3 (Machine-actionable guidance for autonomous
% data retrieval workflows). Requires \usepackage{tcolorbox} in the preamble.
% Drop % Boxed example for Section 2.3 (Machine-actionable guidance for autonomous
% data retrieval workflows). Requires \usepackage{tcolorbox} in the preamble.
% Drop % Boxed example for Section 2.3 (Machine-actionable guidance for autonomous
% data retrieval workflows). Requires \usepackage{tcolorbox} in the preamble.
% Drop \input{section_2_3_example_box} where the box should appear, or paste
% the tcolorbox environment directly into the section.

\begin{tcolorbox}[colback=gray!5,colframe=gray!50,boxrule=0.5pt,
    title={\textbf{Worked example: from per-granule verification to a retrieval recipe}},
    fonttitle=\normalsize]
\small
We have already exercised this workflow end-to-end against a live Hyrax server
(the grep-dap proof-of-concept). The example below recasts that demonstration
for the EDC setting, where collections already carry detailed CMR and CF
metadata and the agent's job is therefore \emph{verification} and
\emph{cross-granule accumulation} rather than recovery from scratch.

\textbf{What the agent does, per collection.} Walking the granules of a
collection through OPeNDAP, the agent (1) confirms that the stated semantic
metadata is internally consistent across granules and consistent with the data
itself, and (2) accumulates collection-level facts that no single granule
record exposes. Each finding is anchored in evidence and graded by confidence,
and the agent explicitly refuses to assert anything it cannot defend.

\begin{itemize}\itemsep1pt
  \item \textbf{Value-range vs.\ declared units.} A small DAP4 hyperslab
  request confirms that a variable declared \texttt{units="kelvin"} actually
  spans $\sim$270--310, not $\sim$$-3$--35. In the proof-of-concept this same
  check caught two genuinely mislabeled variables (a wind speed declared
  \texttt{"mm"}, a pixel count declared \texttt{"C"}).
  \item \textbf{Cross-granule accumulation.} Per-granule spatial/temporal
  bounds, fill-value usage, and valid ranges are aggregated into a single
  collection-level envelope, and divergence between granules is flagged.
  \item \textbf{Cross-checks against CMR.} Accumulated bounds are compared with
  the collection's CMR record; disagreements are surfaced for an expert rather
  than silently overwritten.
  \item \textbf{Retrieval-relevant structure.} DMR++ chunk sizes and layout are
  read directly, yielding concrete subsetting guidance (chunk-aligned request
  shapes, byte ranges, and dataset-specific pitfalls).
\end{itemize}

\textbf{The 2.3 payload.} These verified, evidence-tagged facts are written to
a per-collection \texttt{usage.md} and embedded in the collection's DAP4
responses. A downstream autonomous agent then consumes it in-band, at the
moment of retrieval, and issues a correct, chunk-aligned subsetting request on
the first attempt---without a separate discovery round-trip and without
guessing at units, coordinate orientation, or fill conventions. This is the
difference Part~3 delivers: the curated guidance from Part~2 travels with the
data, turning queries that agents already make into queries that succeed.
\end{tcolorbox}
 where the box should appear, or paste
% the tcolorbox environment directly into the section.

\begin{tcolorbox}[colback=gray!5,colframe=gray!50,boxrule=0.5pt,
    title={\textbf{Worked example: from per-granule verification to a retrieval recipe}},
    fonttitle=\normalsize]
\small
We have already exercised this workflow end-to-end against a live Hyrax server
(the grep-dap proof-of-concept). The example below recasts that demonstration
for the EDC setting, where collections already carry detailed CMR and CF
metadata and the agent's job is therefore \emph{verification} and
\emph{cross-granule accumulation} rather than recovery from scratch.

\textbf{What the agent does, per collection.} Walking the granules of a
collection through OPeNDAP, the agent (1) confirms that the stated semantic
metadata is internally consistent across granules and consistent with the data
itself, and (2) accumulates collection-level facts that no single granule
record exposes. Each finding is anchored in evidence and graded by confidence,
and the agent explicitly refuses to assert anything it cannot defend.

\begin{itemize}\itemsep1pt
  \item \textbf{Value-range vs.\ declared units.} A small DAP4 hyperslab
  request confirms that a variable declared \texttt{units="kelvin"} actually
  spans $\sim$270--310, not $\sim$$-3$--35. In the proof-of-concept this same
  check caught two genuinely mislabeled variables (a wind speed declared
  \texttt{"mm"}, a pixel count declared \texttt{"C"}).
  \item \textbf{Cross-granule accumulation.} Per-granule spatial/temporal
  bounds, fill-value usage, and valid ranges are aggregated into a single
  collection-level envelope, and divergence between granules is flagged.
  \item \textbf{Cross-checks against CMR.} Accumulated bounds are compared with
  the collection's CMR record; disagreements are surfaced for an expert rather
  than silently overwritten.
  \item \textbf{Retrieval-relevant structure.} DMR++ chunk sizes and layout are
  read directly, yielding concrete subsetting guidance (chunk-aligned request
  shapes, byte ranges, and dataset-specific pitfalls).
\end{itemize}

\textbf{The 2.3 payload.} These verified, evidence-tagged facts are written to
a per-collection \texttt{usage.md} and embedded in the collection's DAP4
responses. A downstream autonomous agent then consumes it in-band, at the
moment of retrieval, and issues a correct, chunk-aligned subsetting request on
the first attempt---without a separate discovery round-trip and without
guessing at units, coordinate orientation, or fill conventions. This is the
difference Part~3 delivers: the curated guidance from Part~2 travels with the
data, turning queries that agents already make into queries that succeed.
\end{tcolorbox}
 where the box should appear, or paste
% the tcolorbox environment directly into the section.

\begin{tcolorbox}[colback=gray!5,colframe=gray!50,boxrule=0.5pt,
    title={\textbf{Worked example: from per-granule verification to a retrieval recipe}},
    fonttitle=\normalsize]
\small
We have already exercised this workflow end-to-end against a live Hyrax server
(the grep-dap proof-of-concept). The example below recasts that demonstration
for the EDC setting, where collections already carry detailed CMR and CF
metadata and the agent's job is therefore \emph{verification} and
\emph{cross-granule accumulation} rather than recovery from scratch.

\textbf{What the agent does, per collection.} Walking the granules of a
collection through OPeNDAP, the agent (1) confirms that the stated semantic
metadata is internally consistent across granules and consistent with the data
itself, and (2) accumulates collection-level facts that no single granule
record exposes. Each finding is anchored in evidence and graded by confidence,
and the agent explicitly refuses to assert anything it cannot defend.

\begin{itemize}\itemsep1pt
  \item \textbf{Value-range vs.\ declared units.} A small DAP4 hyperslab
  request confirms that a variable declared \texttt{units="kelvin"} actually
  spans $\sim$270--310, not $\sim$$-3$--35. In the proof-of-concept this same
  check caught two genuinely mislabeled variables (a wind speed declared
  \texttt{"mm"}, a pixel count declared \texttt{"C"}).
  \item \textbf{Cross-granule accumulation.} Per-granule spatial/temporal
  bounds, fill-value usage, and valid ranges are aggregated into a single
  collection-level envelope, and divergence between granules is flagged.
  \item \textbf{Cross-checks against CMR.} Accumulated bounds are compared with
  the collection's CMR record; disagreements are surfaced for an expert rather
  than silently overwritten.
  \item \textbf{Retrieval-relevant structure.} DMR++ chunk sizes and layout are
  read directly, yielding concrete subsetting guidance (chunk-aligned request
  shapes, byte ranges, and dataset-specific pitfalls).
\end{itemize}

\textbf{The 2.3 payload.} These verified, evidence-tagged facts are written to
a per-collection \texttt{usage.md} and embedded in the collection's DAP4
responses. A downstream autonomous agent then consumes it in-band, at the
moment of retrieval, and issues a correct, chunk-aligned subsetting request on
the first attempt---without a separate discovery round-trip and without
guessing at units, coordinate orientation, or fill conventions. This is the
difference Part~3 delivers: the curated guidance from Part~2 travels with the
data, turning queries that agents already make into queries that succeed.
\end{tcolorbox}
 where the box should appear, or paste
% the tcolorbox environment directly into the section.

\begin{tcolorbox}[colback=gray!5,colframe=gray!50,boxrule=0.5pt,
    title={\textbf{Worked example: from per-granule verification to a retrieval recipe}},
    fonttitle=\normalsize]
\small
We have already exercised this workflow end-to-end against a live Hyrax server
(the grep-dap proof-of-concept). The example below recasts that demonstration
for the EDC setting, where collections already carry detailed CMR and CF
metadata and the agent's job is therefore \emph{verification} and
\emph{cross-granule accumulation} rather than recovery from scratch.

\textbf{What the agent does, per collection.} Walking the granules of a
collection through OPeNDAP, the agent (1) confirms that the stated semantic
metadata is internally consistent across granules and consistent with the data
itself, and (2) accumulates collection-level facts that no single granule
record exposes. Each finding is anchored in evidence and graded by confidence,
and the agent explicitly refuses to assert anything it cannot defend.

\begin{itemize}\itemsep1pt
  \item \textbf{Value-range vs.\ declared units.} A small DAP4 hyperslab
  request confirms that a variable declared \texttt{units="kelvin"} actually
  spans $\sim$270--310, not $\sim$$-3$--35. In the proof-of-concept this same
  check caught two genuinely mislabeled variables (a wind speed declared
  \texttt{"mm"}, a pixel count declared \texttt{"C"}).
  \item \textbf{Cross-granule accumulation.} Per-granule spatial/temporal
  bounds, fill-value usage, and valid ranges are aggregated into a single
  collection-level envelope, and divergence between granules is flagged.
  \item \textbf{Cross-checks against CMR.} Accumulated bounds are compared with
  the collection's CMR record; disagreements are surfaced for an expert rather
  than silently overwritten.
  \item \textbf{Retrieval-relevant structure.} DMR++ chunk sizes and layout are
  read directly, yielding concrete subsetting guidance (chunk-aligned request
  shapes, byte ranges, and dataset-specific pitfalls).
\end{itemize}

\textbf{The 2.3 payload.} These verified, evidence-tagged facts are written to
a per-collection \texttt{usage.md} and embedded in the collection's DAP4
responses. A downstream autonomous agent then consumes it in-band, at the
moment of retrieval, and issues a correct, chunk-aligned subsetting request on
the first attempt---without a separate discovery round-trip and without
guessing at units, coordinate orientation, or fill conventions. This is the
difference Part~3 delivers: the curated guidance from Part~2 travels with the
data, turning queries that agents already make into queries that succeed.
\end{tcolorbox}
